\begin{frame}{자주 묻는 질문 2: Block A를 건너뛰면?}
  \textbf{Q. Block A(고전 ML)를 안 배우고 Block B(신경망)로 바로
  넘어가면 안 되나요?}
  \begin{itemize}
    \item 이론적으로는 되지만, \textbf{권장하지 않는다}.
    \item Week 6 정규화, Week 7 트리 앙상블의 \textbf{``편향-분산
      트레이드오프''}와 ``train/val/test 분리''는 신경망에도
      \textbf{그대로 적용} --- Week 9의 드롭아웃·배치정규화는 정확히
      같은 원리.
    \item 지금도 정형 데이터에서는 고전 ML이 신경망을 이기는 경우가
      흔해서 \textbf{실무에서 바로 쓸모}가 있다.
  \end{itemize}
  \vskip0.6em
  \textbf{Q. 수학이 약한데 따라갈 수 있나요?}
  \begin{itemize}
    \item 1.3절의 세 가지 수학(내적, 연쇄법칙, 베이즈 정리)이
      \textbf{사실상 전부}. 그 이상의 새 수학 도구는 거의 등장하지 않고,
    \item 대신 그 세 가지를 \textbf{점점 더 복잡한 상황에 반복 적용}하는
      것이 이번 학기의 실제 난이도. 낯선 표기가 나올 때마다 1.3절을
      참고하는 습관이면 충분.
  \end{itemize}
\end{frame}
